\documentclass[12pt]{article}

\usepackage[a4paper, total={6in, 8in}]{geometry}
\usepackage{textcomp}

\begin{document}
\setlength{\parindent}{0pt}

\def \t {\textrightarrow}
\def \v {\vspace{1em}}
\def \bi {\begin{itemize}}
\def \ei {\end{itemize}}
\def \s[#1] {\section*{#1}}
\def \ss[#1] {\subsection*{#1}}
\def \sss[#1] {\subsubsection*{#1}}

\s[Metabolismo]
Immagazziniamo l'energia nell atp \t la prendiamo dai nutrienti \\
Il metabolismo autorotrofo \t organismi che sono in grado di produrre molecole da se, come le piante ed alcuni batteri \\
Il metabolismo catabolico = degradazione delle molecole per produrre energia \\
Anabolismo = parte del metabolismo che produce molecole \\
Sei io produco energia parlo di una reazione che ha bilancio esoergonico \t mentre processo anabolico ha reazione con bilancio positivo (perche immagazzino energia, che pero viene sottratta dall'ambiente) \\
Differenza tra organismi autorotrofi ed eterotrofi:
\bi
  \item autorotrofi usano energia (per esempio solare) e sintetizzano le molecole complesse
  \item eterotrofi prendono ??
\ei
La via metabolica porta ad usare il glucosio come fonte di energia \t le reazioni nella gran parte sono redox (RIVEDERLE) \\
Cicli in cui c'è passaggio di elettroni \t per farlo ho bisogni di molecole apposite, e quindi i conenzimi (molecole che attivano gli enzimi) \\
Le molecole che portano alle reazioni metaboliche con enzimi \t continuano a passare continuamente da condizione ridotta e ossidata (es. NAD) \\
Chiedere da qui in po 
\end{document}
